\documentclass[12pt]{article}
\usepackage[utf8]{inputenc}
\usepackage[T1]{fontenc}
\usepackage{hyperref}
\title{Wadaginam Language Resources for K-12 Curriculum Development}
\author{Compiled by Claude Sonnet 5 (claude-sonnet-5) for Andr\'as Kornai}
\date{}
\begin{document}
\maketitle

\section{Introduction}
Wadaginam (ISO 639-3 wdg; Glottocode wada1263) is a divergent Madang language of the Adelbert Range in Madang Province, Papua New Guinea. It is spoken south of Bogia. Wikipedia gives about 950 native speakers (2003). Joshua Project lists about 1,800 people, with low literacy. Neighbouring groups are Tanguat and Inapang.

\section{Primary Educational Resources}
No primers, textbooks or school programmes for Wadaginam were found in search results.

\section{Linguistic and Reference Materials}
A wordlist drawn from Z'graggen (1980) is hosted at the Trans New Guinea database: \url{https://transnewguinea.org/language/wadaginam?page=3}. Wikipedia gives an overview: \url{https://en.wikipedia.org/wiki/Wadaginam_language}. No full grammar or dictionary was found.

\section{Religious and Community Resources}
Joshua Project reports no Bible portions, no Jesus Film and no audio Bible: \url{https://joshuaproject.net/languages/wdg} and \url{https://joshuaproject.net/people_groups/15814}. Scripture Earth has a page for this code, but the search results showed no content in the language: \url{https://www.scriptureearth.org/00eng.php?iso=wdg}.

\section{Digital and Audio Resources}
No audio or video resources were found in search results.

\section{Curriculum Development Context}
Documentation is thin, and a curriculum would depend on new orthography work and community partnership. Tok Pisin and English are the likely languages of wider schooling.
\end{document}